% ==============================================================
%  Chapter 10: Top 10 Concepts You Must Understand
% ==============================================================
\chapter{Top 10 Concepts You Must Understand}
\label{ch:concepts}

These are the ten concepts that separate people who use AI effectively from people who
are constantly frustrated by it.
You do not need to understand the mathematics behind them, but you need to understand
what they \emph{mean in practice}.

% --------------------------------------------------------------
\section{Concept 1: Prompt Engineering}
\label{sec:prompt-engineering}
% --------------------------------------------------------------

\begin{conceptbox}[Concept 1: Prompt Engineering]
The single most important skill for getting value from AI.
A better prompt almost always produces a better result than a more powerful model.
\end{conceptbox}

\textbf{What it means:} Prompt engineering is the practice of structuring your input to
an AI model to get the most useful output.
It is not magic or trick-taking --- it is clear communication.

\textbf{The core principles:}
\begin{itemize}
  \item \textbf{Role:} Tell the model who it should act as.
    \textit{``You are a senior Python developer performing a security code review.''} \\
    This activates the relevant knowledge domain and tone.
  \item \textbf{Context:} Provide the background needed to give a useful answer.
  \item \textbf{Task:} State the goal precisely.
  \item \textbf{Format:} State the output format (bullet list, table, code block,
    one paragraph).
  \item \textbf{Constraints:} Say what to avoid, what length to target, what style
    to use.
\end{itemize}

\textbf{The difference it makes:}
\begin{center}
\begin{tabular}{ll}
\toprule
\textbf{Vague prompt} & \textbf{Engineered prompt} \\
\midrule
``Fix my code'' & ``Review this Python function for correctness, edge cases, and \\
& readability. Return: (1) a list of issues by severity, \\
& (2) a corrected version, (3) explanation of each change.'' \\
\bottomrule
\end{tabular}
\end{center}

% --------------------------------------------------------------
\section{Concept 2: Context Windows and Memory}
\label{sec:context-window}
% --------------------------------------------------------------

\begin{conceptbox}[Concept 2: Context Windows and Memory]
AI has no persistent memory between conversations.
Within a conversation, it can only see what fits in its context window.
\end{conceptbox}

\textbf{What it means:}
A context window is the maximum amount of text (measured in tokens) that the model
can ``see'' at once --- your message, its previous responses, and anything you paste.
Models have context limits ranging from 8,000 to 200,000+ tokens depending on the
model.

\textbf{Practical implications:}
\begin{itemize}
  \item Long conversations or large file pastes eventually push early context out of
        the window --- the model effectively forgets what was said at the start.
  \item There is no persistent memory between sessions unless a tool explicitly saves
        and re-injects it.
  \item For large codebases, tools like Cursor and Copilot use selective retrieval
        (embeddings and search) to bring relevant context in, rather than loading
        everything at once.
\end{itemize}

\textbf{What to do:}
Restate important context at the top of each new session.
Don't assume AI remembers your tech stack, project constraints, or decisions from a
previous conversation.

% --------------------------------------------------------------
\section{Concept 3: Hallucination}
\label{sec:hallucination}
% --------------------------------------------------------------

\begin{conceptbox}[Concept 3: Hallucination]
AI generates the most statistically likely continuation of your prompt.
Sometimes that continuation is false, but stated with complete confidence.
\end{conceptbox}

\textbf{What it means:}
Hallucination is when AI generates incorrect information --- a non-existent function,
a wrong date, a fabricated citation, a plausible but false code pattern --- without
flagging any uncertainty.
The problem is not that AI is sometimes wrong.
The problem is that it is wrong and confident in the same way it is right and confident.

\textbf{Where hallucination is most common:}
\begin{itemize}
  \item Specific version numbers, release dates, and API signatures
  \item Citations, paper titles, and author names
  \item Specific statistics and numerical facts without sources
  \item Rare or niche topics with less training data
\end{itemize}

\textbf{How to mitigate it:}
\begin{itemize}
  \item Verify all specific facts against primary sources
  \item Test all code before trusting it
  \item Ask: \textit{``How confident are you in this and how would I verify it?''}
  \item Use search AI (Perplexity) when you need cited, verifiable facts
\end{itemize}

% --------------------------------------------------------------
\section{Concept 4: RAG (Retrieval-Augmented Generation)}
\label{sec:rag}
% --------------------------------------------------------------

\begin{conceptbox}[Concept 4: RAG]
A technique for giving AI access to your specific documents and data without retraining
the model.
\end{conceptbox}

\textbf{What it means:}
RAG works in two steps:
\begin{enumerate}
  \item Your documents are converted to numerical vectors (embeddings) and stored in
        a vector database.
  \item When you ask a question, the most relevant document chunks are retrieved and
        injected into the model's context alongside your question.
\end{enumerate}
This allows AI to answer questions about your internal knowledge base, documentation,
or data without the model having been trained on that data.

\textbf{Where you see it:}
Notion AI, Confluence AI, GitHub Copilot's codebase understanding, and any
``chat with your documents'' tool is using some form of RAG.

\textbf{Why it matters for professionals:}
RAG is how you deploy AI over proprietary, private, or up-to-date information.
Understanding this concept helps you evaluate knowledge-AI tools and assess their
limitations (retrieval quality and context window constraints).

% --------------------------------------------------------------
\section{Concept 5: AI Agents and Tool Use}
\label{sec:agents}
% --------------------------------------------------------------

\begin{conceptbox}[Concept 5: Agents and Tool Use]
An AI agent takes sequences of actions autonomously: calling tools, writing code,
reading files, and executing steps --- not just generating text.
\end{conceptbox}

\textbf{What it means:}
An agent receives a high-level goal, plans a sequence of steps, calls tools (a web
browser, a code interpreter, a file system, an API), evaluates the result, and
continues until the goal is met.
This is fundamentally different from a chat model that only generates text.

\textbf{Where you see it:}
Cursor Agent, GitHub Copilot Workspace, Devin, and n8n AI nodes are examples.

\textbf{Practical implications:}
\begin{itemize}
  \item Agents can make mistakes that are hard to reverse (deleting files, sending
        emails, modifying a database). Review before allowing irreversible actions.
  \item The goal specification matters enormously --- unclear goals lead to
        incorrect sequences of steps.
  \item Current agents are good at well-defined, bounded tasks.
        Open-ended or ambiguous tasks still fail often.
\end{itemize}

% --------------------------------------------------------------
\section{Concept 6: Temperature and Sampling}
\label{sec:temperature}
% --------------------------------------------------------------

\begin{conceptbox}[Concept 6: Temperature]
Temperature controls how creative vs.\ deterministic the model's outputs are.
\end{conceptbox}

\textbf{What it means:}
At low temperature (near 0), the model picks the most likely word at each step ---
outputs are consistent and conservative.
At high temperature (near 1), more surprising word choices are possible --- outputs
are more creative and varied, but also less reliable.

\textbf{Practical implications:}
\begin{itemize}
  \item For code generation, debugging, and factual Q\&A: use low temperature (0.0--0.2)
  \item For brainstorming, creative writing, and generating diverse options: use higher
        temperature (0.5--0.8)
  \item Most chat interfaces set temperature for you; but if using an API directly,
        this is the first setting to configure
\end{itemize}

% --------------------------------------------------------------
\section{Concept 7: Fine-Tuning vs.\ Prompting}
\label{sec:finetuning}
% --------------------------------------------------------------

\begin{conceptbox}[Concept 7: Fine-tuning vs.\ Prompting]
Prompting instructs a general model. Fine-tuning trains a new model on your data.
Use prompting first; fine-tune only when prompting hits a real wall.
\end{conceptbox}

\textbf{What it means:}
\begin{itemize}
  \item \textbf{Prompting:} Instructing a general-purpose model with a text prompt.
    No training required, immediately testable, adjustable at any time.
  \item \textbf{Fine-tuning:} Training a new model variant on your domain-specific data.
    Expensive, requires curated data, and produces a model that is better at your
    specific task but less flexible.
\end{itemize}

\textbf{When to fine-tune:}
Fine-tuning makes sense when you have thousands of high-quality examples, your task
is highly specialised and repetitive (consistent format, domain-specific terminology),
and prompting has demonstrably hit its limits.

\textbf{What to do instead:} RAG covers most knowledge retrieval needs.
Few-shot prompting (giving 3--5 examples in your prompt) covers most style and format
needs.
Fine-tuning is usually a last resort, not a first step.

% --------------------------------------------------------------
\section{Concept 8: Embeddings and Semantic Search}
\label{sec:embeddings}
% --------------------------------------------------------------

\begin{conceptbox}[Concept 8: Embeddings]
Embeddings are numerical representations of meaning that allow AI to find
conceptually similar content, not just keyword matches.
\end{conceptbox}

\textbf{What it means:}
An embedding converts text into a vector of numbers that encodes its meaning.
Texts with similar meaning have similar vectors, regardless of the words used.
This enables semantic search: finding documents about ``container orchestration'' even if
the document says ``Kubernetes workload scheduling''.

\textbf{Why it matters:}
Embeddings underpin RAG, code search in Copilot, semantic similarity features, and
recommendation systems.
If you are building or evaluating any AI-powered search or knowledge retrieval tool,
embedding quality is the key variable.

% --------------------------------------------------------------
\section{Concept 9: The Token Economy}
\label{sec:tokens}
% --------------------------------------------------------------

\begin{conceptbox}[Concept 9: Tokens and Cost]
AI models process text as tokens (roughly 0.75 words each).
Both input (your prompt) and output (the response) cost tokens.
Cost, latency, and context limits all depend on token count.
\end{conceptbox}

\textbf{What it means:}
\begin{itemize}
  \item Longer prompts cost more to process and count against the context window
  \item API pricing is usually per 1,000 input tokens and per 1,000 output tokens
    (output costs more)
  \item A 100-page document is roughly 75,000 tokens --- expensive to process in
    full at every query; RAG is the solution
\end{itemize}

\textbf{Cost management:}
\begin{itemize}
  \item Be precise in prompts --- not because AI responds better to brevity,
    but because concise prompts cost less
  \item Use smaller/cheaper models for simple classification or extraction tasks;
    reserve larger models for complex reasoning
  \item Cache repeated system prompts in API implementations
\end{itemize}

% --------------------------------------------------------------
\section{Concept 10: The Human-in-the-Loop Principle}
\label{sec:hitl}
% --------------------------------------------------------------

\begin{conceptbox}[Concept 10: Human-in-the-Loop]
AI generates; humans decide.
The human in the loop is not a bottleneck --- they are the quality gate, the
accountability holder, and the source of context AI cannot have.
\end{conceptbox}

\textbf{What it means:}
Human-in-the-loop (HITL) is the principle that humans remain in the decision-making
chain for any consequential AI output.
The degree of human involvement varies by risk level:
\begin{itemize}
  \item \textbf{High risk} (production code, customer communication, financial data):
    human reviews every output before use
  \item \textbf{Medium risk} (internal drafts, analysis summaries): human reviews
    spot-checks and retains override authority
  \item \textbf{Low risk} (formatting, classification of internal data): automation
    acceptable with monitoring
\end{itemize}

\textbf{Why it matters:}
As AI becomes more capable, the temptation to remove humans from the loop grows.
But AI failures in high-stakes domains are costly and hard to detect automatically.
The skill is calibrating the right level of oversight for the right risk --- not
applying maximum human review everywhere (too slow) or minimum review everywhere
(too risky).
